\documentclass[25pt, a0paper, landscape, innermargin=12mm, blockverticalspace=10pt, colspace=12pt, subcolspace=0pt]{tikzposter}

\usepackage[utf8]{inputenc}
\usepackage[T1]{fontenc}
\usepackage{polski}
\usepackage{babel}
\usepackage{booktabs}
\usepackage{graphicx}
\usepackage{qrcode}
\usepackage{xcolor}
\usepackage{amsmath}
\usepackage{enumitem}
\usepackage{ragged2e}
\usepackage{multicol}

\definecolor{AccentBlue}{HTML}{1A3A5C}
\definecolor{LightBlue}{HTML}{EBF2F8}
\definecolor{MediumBlue}{HTML}{3A6B9F}
\definecolor{DarkText}{HTML}{1A1A1A}

\usetheme{Default}
\definecolorstyle{flatBlue}{
    \colorlet{colorOne}{white}
    \colorlet{colorTwo}{LightBlue}
    \colorlet{colorThree}{MediumBlue}
}{
    \colorlet{backgroundcolor}{white}
    \colorlet{framecolor}{white}
    \colorlet{titlefgcolor}{white}
    \colorlet{titlebgcolor}{AccentBlue}
    \colorlet{blocktitlefgcolor}{white}
    \colorlet{blocktitlebgcolor}{AccentBlue}
    \colorlet{blockbodyfgcolor}{DarkText}
    \colorlet{blockbodybgcolor}{white}
    \colorlet{innerblockbodybgcolor}{LightBlue}
    \colorlet{innerblockbodyfgcolor}{DarkText}
    \colorlet{innerblocktitlebgcolor}{MediumBlue}
    \colorlet{innerblocktitlefgcolor}{white}
    \colorlet{notebgcolor}{white}
    \colorlet{notefgcolor}{DarkText}
    \colorlet{notefrcolor}{MediumBlue}
}
\usecolorstyle{flatBlue}

\usetitlestyle{Default}
\settitle{
    \centering
    \vbox{
        \color{white}\bfseries\fontsize{42}{50}\selectfont
        Porównanie metod transfer learning i niestandardowej sieci CNN\\[4mm]
        \fontsize{26}{32}\selectfont\color{white!90}
        w zadaniu klasyfikacji emocji na podstawie ekspresji twarzy\\[6mm]
        \fontsize{20}{24}\selectfont\color{white!80}
        Autorzy: Karol Koenig, Wiktoria Sikora, Nikola Niełacna
    }
}
\usebackgroundstyle{Default}
\usenotestyle{Default}

\setlist{nosep,leftmargin=*,itemsep=3pt,topsep=3pt,parsep=0pt}

\begin{document}

\maketitle[width=0.985\textwidth,titletoblockverticalspace=8mm]

\begin{columns}

% ===== LEWA KOLUMNA (28%) =====
\column{0.28}

\block{Wprowadzenie}{
    \justify
    Rozpoznawanie ekspresji twarzy (\textit{Facial Expression Recognition}, FER)
    jest zadaniem z~podobszaru \textbf{computer vision}, polegającym na klasyfikacji
    obrazów twarzy do jednej z~podstawowych kategorii emocji.
    W~projekcie wykorzystano zbiór danych \textbf{Kaggle Emotion Detection FER},
    zawierający \textbf{35\,887 obrazów} w~skali szarości o~rozdzielczości
    $48\times48$ pikseli, z~podziałem na 7 klas emocji:
    \begin{center}
        \textit{angry, disgusted, fearful, happy, neutral, sad, surprised}
    \end{center}
    \bigskip
    \small
    \textbf{Literatura:}\newline
    [1] M. Xu et al., \textit{A Comprehensive Survey of Image Augmentation Techniques
    for Deep Learning}, Pattern Recognition 137, 2023.\newline
    [2] K. He et al., \textit{Deep Residual Learning for Image Recognition},
    IEEE CVPR, 2016.\newline
    [3] K. Simonyan, A. Zisserman, \textit{Very Deep Convolutional Networks for
    Large-Scale Image Recognition}, ICLR, 2015.
    \normalsize
}

\block{Problem badawczy}{
    \justify
    \begin{enumerate}
        \item \textbf{Porównanie} trzech architektur \textit{transfer learning}
              (ResNet18, VGG16, ResNet50) z~autorską siecią \textbf{CustomCNN}
              trenowaną od podstaw.
        \item Ocena skuteczności \textbf{ważenia klas} w~funkcji kosztu
              jako metody radzenia sobie z~problemem niezbalansowanych danych.
        \item Identyfikacja \textbf{najczęściej mylonych par emocji}
              i~analiza przyczyn błędów klasyfikacji.
    \end{enumerate}
}

\block{Metryki per-class}{
    \centering
    \includegraphics[width=\linewidth]{img/per_class_metrics.pdf}\\[2mm]
    {\small Rysunek 3. Precision, Recall i F1 dla każdej klasy i modelu.}
}

% ===== KOLUMNA CENTRALNA (42%) =====
\column{0.42}

\colorlet{saveTitleBg}{blocktitlebgcolor}
\colorlet{saveBodyBg}{blockbodybgcolor}
\colorlet{blocktitlebgcolor}{AccentBlue}
\colorlet{blockbodybgcolor}{LightBlue}

\block{\centering Porównanie modeli --- średnie metryki (10-fold CV)}{
    \centering
    \includegraphics[width=\linewidth]{img/model_comparison.pdf}\\[2mm]
    {\small Rysunek 1. Porównanie średnich metryk dla 4 modeli.}

    \bigskip
    \begin{center}
    {\renewcommand{\arraystretch}{1.3}
    \begin{tabular}{lcccc}
        \toprule
        \textbf{Model} & \textbf{Accuracy} & \textbf{Balanced Acc.} & \textbf{Precision} & \textbf{Recall} \\
        \midrule
        VGG16      & 0,6137(59) & \textbf{0,5915(86)} & \textbf{0,6215(80)} & \textbf{0,6137(59)} \\
        CustomCNN  & 0,5881(47) & 0,5627(77)           & 0,5853(57)           & 0,5881(47) \\
        ResNet18   & 0,5814(58) & 0,5517(68)           & 0,5882(89)           & 0,5814(58) \\
        ResNet50   & 0,5381(38) & 0,5012(61)           & 0,5423(71)           & 0,5381(38) \\
        \bottomrule
    \end{tabular}}\\[2mm]
    {\small Tabela 1. Średnie wartości metryk (10-fold CV) z~odchyleniem standardowym.}
    \end{center}

    \bigskip
    \justify
    \textbf{VGG16} osiągnął najwyższą zbalansowaną dokładność
    (Balanced Acc $\approx 0{,}59$). \textbf{CustomCNN} ($\approx 0{,}56$),
    trenowana 60 epok od zera, przewyższa oba modele ResNet. ResNet50 ---
    najgłębsza architektura --- osiąga najsłabsze wyniki, co sugeruje
    trudności z~efektywnym fine-tuningiem przy ograniczonej liczbie 
    odmrożonych warstw, co było wymuszone ograniczeniami sprzętowymi.
}

\colorlet{blocktitlebgcolor}{saveTitleBg}
\colorlet{blockbodybgcolor}{saveBodyBg}

\block{Macierze pomyłek}{
    \centering
    \includegraphics[width=\linewidth]{img/confusion_matrices.pdf}\\[2mm]
    {\small Rysunek 2. Znormalizowane macierze pomyłek dla 4 modeli.}

    \bigskip
    \justify
    \textbf{Najczęstsze błędy (wszystkie modele):}
    \textit{disgusted $\to$ angry} (17--23\%), \textit{fearful $\leftrightarrow$ sad}
    (16--19\%), \textit{neutral $\to$ sad} (15--21\%). Emocje \textit{happy}
    (F1 $\approx$ 0{,}81) i~\textit{surprised} (F1 $\approx$ 0{,}73) rozpoznawane
    są najpewniej. \textit{Fearful} (F1 $\approx$ 0{,}46) i~\textit{disgusted}
    (F1 $\approx$ 0{,}49) pozostają najtrudniejsze. Pomyłki fearful $\leftrightarrow$ sad
    oraz neutral $\leftrightarrow$ sad wskazują na wizualne podobieństwo tych emocji
    w~zbiorze danych.
}

% ===== PRAWA KOLUMNA (30%) =====
\column{0.30}

\block{Protokół eksperymentalny}{
    \justify

    \textbf{Walidacja:} RepeatedStratifiedKFold ($n_{\mathrm{splits}}{=}2$,
    $n_{\mathrm{repeats}}{=}5$). Zbiory treningowy i~testowy scalone
    przed podziałem (35\,887 obrazów łącznie).

    \textbf{Transfer learning:} modele ImageNet-pretrained (ResNet18, VGG16, ResNet50).
    Odblokowane ostatnie warstwy konwolucyjne +~wymieniona głowa klasyfikatora.
    Fine-tuning: 10 epok/fold. Dwie wartości learning rate:
    $10^{-4}$ dla warstw pretrenowanych, $10^{-3}$ dla nowej głowy.
    Przetwarzanie danych: resize + crop + normalizacja.

    \textbf{CustomCNN:} architektura 4-blokowa (Conv $\to$ ReLU $\to$ BatchNorm
    $\to$ MaxPool $\to$ Dropout2D), trenowana \textbf{60 epok} od zera.
    Adam (lr $= 10^{-4}$, weight\_decay $= 10^{-4}$). Wejście: $1 \times 48 \times 48$ (grayscale).
    Przetwarzanie danych: resize + crop + normalizacja.

    \textbf{Funkcja kosztu:} CrossEntropyLoss z~wagami klas liczonymi dynamicznie
    dla każdego foldu: $w_c = N_{\mathrm{train}} / (K \cdot N_c)$, gdzie $K = 7$,
    $N_c$ --- liczba próbek klasy $c$ w~podzbiorze treningowym foldu.

    \textbf{Metryki:} Accuracy, Balanced Accuracy, Weighted Precision,
    Weighted Recall, Confusion Matrix.

    \textbf{Sprzęt:} AMD Radeon RX 6700 XT, DirectML, PyTorch 2.4.
}

\block{Wnioski}{
    \justify
    \begin{enumerate}
        \item \textbf{VGG16} (Balanced Acc $\approx 0{,}59$) osiąga najlepszą
              zbalansowaną dokładność spośród testowanych modeli --- ta
              architektura lepiej sprawdza się w zadaniu.
        \item \textbf{CustomCNN} (Balanced Acc $\approx 0{,}56$), trenowana
              60 epok od zera, przewyższa ResNet18 i~ResNet50. Mniejsza ilość parametrów
              pozwala na bardziej dogłębne uczenie. Wynik potwierdza
              wartość dedykowanej architektury dla konkretnych zastosowań.
        \item Ważenie klas w~funkcji kosztu \textbf{częściowo łagodzi}
              problem niezbalansowania we wszystkich modelach, ale klasa
              \textit{disgusted} (1{,}5\% danych) pozostaje wyzwaniem.
              Rozwiązaniem może być zastosowanie technik generatywnych (GAN, Diffusion).
        \item Pomyłki \textbf{fearful $\leftrightarrow$ sad} oraz
              \textbf{neutral $\leftrightarrow$ sad} wskazują na wizualne
              podobieństwo tych emocji.
    \end{enumerate}
    \bigskip
}

\block{Kod źródłowy}{
    \centering
    \qrcode[height=2.5cm]{https://github.com/JestemSikora/pytorch-emotion-detector}\\[2mm]
    {\small Repozytorium projektu}
}

\end{columns}

\end{document}
